%! Zeros \& Graphs of Polynomials — Unit 3, Lesson 3.6 — Student Packet Cover Sheet
% UNITS 3+ PACKET-FIRST shape (LESSON_SHAPE.md §0a, 2026-11-14). 12pt; \coverbanner sizes the
% band. THREE rows (no extra credit, 2026-11-17) (a contents list, no Score column): Warm-Up, Guided Notes & Practice, the
% DeltaMath homework (online, not in this packet). Homework
% is due the first class after two study halls — course policy of record — never "next class".
% Keep in Mind matches the Closing Practice's Watch-out line on the deck.
\documentclass[12pt]{article}
\usepackage{algebra2-article}
\usepackage{algebra2-boxes}

\begin{document}

\coverbanner{Unit 3 \quad Polynomial Functions}{Lesson 3.6 \quad Zeros \& Graphs of Polynomials}

\namedateperiod
\vspace{0.18in}

% Learning targets use the lesson's FORMAL VOCABULARY in bold — the definition frames name it.
\boxguard[14]
\begin{learningtargetbox}
\small
\begin{enumerate}[leftmargin=1.5em, itemsep=5pt]
  \item \ldots{} use the \textbf{Fundamental Theorem of Algebra} to say how many
        \textbf{complex zeros} a polynomial has (its degree, counted with
        \textbf{multiplicity}), and use the degree and a graph's $x$-intercepts to find the
        \textbf{number and type} of zeros --- real or imaginary.
  \item \ldots{} explain why an imaginary zero $a+bi$ of a polynomial with real coefficients
        always comes with its \textbf{complex conjugate} $a-bi$.
  \item \ldots{} write a polynomial from its zeros in \textbf{factored form}
        $a(x-r_1)(x-r_2)\cdots$, supply a missing conjugate, expand it to
        \textbf{standard form}, and find $a$ from one more point.
  \item \ldots{} match a factored polynomial to its graph and defend the match with its
        \textbf{end behavior}, whether it \textbf{crosses} or \textbf{touches} at each zero,
        and its \textbf{$y$-intercept} --- and write a possible equation from a graph.
\end{enumerate}
\end{learningtargetbox}

\vspace{0.14in}

\boxguard[14]
\begin{tocbox}
\small
\renewcommand{\arraystretch}{1.5}
\begin{tabularx}{\linewidth}{@{} c l X @{}}
\rowcolor{forestbg}
\textbf{\#} & \textbf{Component} & \textbf{Description} \\
\midrule
1 & Warm-Up & Multiplying binomials and a conjugate pair, zeros and multiplicity, and every
              solution of $(x-2)(x^2+9)=0$ \\
2 & Guided Notes \& Practice & Three rounds --- counting zeros with the Fundamental Theorem,
              writing a polynomial from its zeros, and matching equations to graphs --- then the
              Closing Practice, scored out of 4 \\
3 & Homework & Assigned on DeltaMath: \textbf{A2: Lesson 3.6: Zeros \& Graphs of Polynomials} --- complete it
              online; it is not in this packet. \textbf{Scored 1 to 5, due the first class after
              two study halls.} \\
\end{tabularx}
\end{tocbox}

\vspace{0.14in}

% Keep in Mind is a CONTENT summary — the lesson's definitions and the distinction it turns on.
\boxguard[8]
\begin{remindbox}
\small
A polynomial of degree $n$ has exactly $n$ \textbf{complex zeros}, counted with
\textbf{multiplicity} (\textbf{Fundamental Theorem of Algebra}). With real coefficients, an
imaginary zero $a+bi$ always brings its \textbf{conjugate} $a-bi$, so imaginary zeros come in
pairs; the real zeros are the $x$-intercepts, and the rest are imaginary. Each zero $r$ gives the
factor $(x-r)$: a zero at $x=3$ is the factor $(x-3)$, \emph{not} $(x+3)$. At an odd power the
graph \textbf{crosses} the $x$-axis; at an even power it \textbf{touches} and turns back --- so a
touch is \emph{squared}, never a single factor.
\end{remindbox}

\end{document}
